\section{Related Work}
\subsection{Alternative Approaches for Encoding PM-Reg}
We consider several alternative approaches for encoding
PM-Reg. The first choice is to use a general proof system
with the support of lookup-arguments.

HyperPlonk \cite{HyperPlonk} improves Plonk by replacing
the univariate polynomial quotient check with multi-variate
polynomial commitment and evaluation proof using sumchecks.
This increases the proof size to $\log(n)$, however,
removes dependence on FFT (thus improving parallelism)
and allows it to support high degree customized gates.
However, like PlonkUp \cite{PlonkUp}, it implements the
Plookup protocl \cite{plookup} which has prover
complexity depending on the lookup table size.
Another restriction of \cite{plookup} and \cite{PlonkUp}
is that the lookup table size cannot exceed
the circuit size, which in our case is a much larger value.

\subsection{Folding and Proof Recursion}

Detailed comments: {\bf Accumulation:}
Halo \cite{Halo} uses cyclic curves (but without pairing)
to achieve low cost proof recursion (however it's recursive
overhead is $\log(n)$).
PCD \cite{PCD20}: It presentes the accumulator scheme to realize
IVC (incremental verification). The basic idea is that an accumulator
is an object of aggregated statement and proof, whose size
does not increase with the number of (statement, proof) that
are aggregated. A new accumulator can be generated from 
an old accumulator, and it recursively proves that all prior
instances are correct. A final decider (which has the same
verifier cost of verifying a single circuit) decides the
validity of all $n$ copies of the circuit.  Compared with Nova (Folding)
\cite{Nova}, it does not use an extended R1CS system, but does
use a similar error term to cancel out the cross terms when
accumulation is applied. It commits to $A \cdot z$, $B \cdot z$,
$C \dot z$ and has a slightly higher concrete cost than Nova.
ProtoStar \cite{ProtoStar} further extends the accmulation scheme,
it accumulates/folds a running accumulator with a NARK (not necessarily)
proof. It accomplishes the non-uniform circuit as SuperNova,
and it provides better verifier circuit complexity using
Hab22 \cite{Hab22} technique ($O(1)$ verifier complexity other than
$O(|t| \log(|T|)$ of HyperNova.  Note that its verifier circuit
has only one more hash operation for the lookup (and others 
error terms computation have the same complexity of others, and
embedded with other constraints). {\bf However, notice that
the cost of $|T|$ has to be paid in the last decider!} (It looks like
if it sacrifices the no trusted-up, it can still be independent
of table size). Note that due to extra long vectors of
$\vec{m}$, these constraints cannot be incorporated into the standard
columns of plonk system. For an extra of $6|t|$ group operations
is needed on the prover side for each step of recursion of
lookup tables, in addition to the standard cost of encoding columns.
Origami \cite{Origami} concurrently proposes a similar folding
skeme as ProtoStar using error terms for AIR, and it applies to
plookup (actually Halo2 lookup) by error terms of polynomial maps.
Note that plookup table has to be the same as number of constraints (cannot
handle large table). Its passing of randoms input is subsumed
by commit-and-prove feature of halo2 (multi-stages) anyway.
DeepThought \cite{DeepThought} further instantiate GKR protocol
over the ProtoStar framework. It makes two contributions: (1) GKR
allows to commit to input and ouput wires only (thus not committing
to all wires); The paper modifies GKR to bivariate to lower
its cost in the ProtoStar framework (reducing rounds)
(2) adapting the Hab22 technique, it can prove
a {\em dynamic} memory table read/write operations (by permuting
read/write records and also add addr/index into Hab22 formula), thus
achieving constant recursion cost per mem op, just like
lookup table in ProtoStar. {\em Altnertive solution: maybe
permutation of memory transcript is good enough. Then for
initial memory read and final memory write in a sorted list they
all can be performed using CQ protocol anyway.}
ProtoGalaxy \cite{ProtoGalaxy} further cuts the recursion cost
of extra $3$ group operations from ProtoStar by the following trick:
it folds multiple instances, and it uses a Lagrange polynomial to encode
the error terms. So instead of folding instances using a random
factor, it allows verifier to use a new random for the Lag based poly
encoding all instances. This cuts the group operations (replacing
with a log number of field operations) in the ``marginal cost" (but
note that linear combination of group elements of claims still
have to be performed in the verifier circuit).
In a very recent work \cite{AccCode}, the accumulation scheme
is achieved by symmetric crypto pritmitices. The basic idea is
to replace the Pedersen commitment of witness vectors by a Merkle
tree of an RS-code of witness. Then the homomorphic group operation
is replaced by a protocol of verify code validity. This moves away
from pub-key assumptions. For decider, its complexity is $n/\log(n)$,
also we can avoid non-native field arithmetics and cycle curve,
if we encode it in BN254?
\cite{ARC} furtuer improves \cite{AccCode} by preserving distance
of RS-code.
Recently Mira \cite{Mira} extends ProtoStar by allowing
group elements in the Special-Sound Protocol and
simulate the running of pairing in SSP, and shows
better performance than Groth16 proof packing ({\em However,
not sure how the discrete logarithm of Groth16 proofs were
obtained? Authors never responded}).


{\bf Folding}
Nova \cite{Nova}: it presents the {\em folding} technique which folds
two relaxed R1CS instance into one. The folding prover at each step
pays two MSMs, the folding verifier pays two group scalar multiplications.
When applied to IVC, the verifier circuit is embedded in the IVC
proof recursively. The proof size is linear, but using a cyclic curve
with sumcheck, the proof size can be compressed to log of circuit size
using a variation of Spartan. 
Sangria \cite{Sangria} implements the same idea of NOVA by
folding PLONK relations. Note that it is subsumed by
ProtoStar.
In \cite{NovaCycle}: a rigorous proof is given
for the correctness of Nova when implemented over
a cycle of curves and correcting one vulnerability in its
implementation.
SuperNova  \cite{SuperNova} further
improves Nova to handle non-uniform R1CS instances, which allows the prover
to pay cost only on instructions that is being executed.
Given $\ell$ circuits, the idea is to maintain maintain
$\ell$ instances of R1CS instances, but when applying
folding, only fold the actual circuit being executed.
Note that the IVC proof size is still linear to the total
of the $\ell$ circuit size, also when proving the circuit
the prover has to maintain the memory for all $\ell$ circuits.
HyperNova \cite{HyperNova} further extends the folding idea
to CCS (a combination of R1CS, Plonky, and AIR specification)
which supports high digree gates. The idea is to model the system
using multi-linear polynomials whose evaluation can be
specified using sumcheck techniques, which results in verifier circuit
of log size of field ops plus one scalar multiplication. The interesting
part of this is that the more expressive the spec system, the easier
the folding? The IVC proof size is still linear, and needs a cyclic
curve for succinct proof. 
CycleFold \cite{CycleFold} further
improves the Nova scheme by not running a dual R1CS system 
over a cycle of curves, but simply outsourcing a few
non-native field and group operations to a second cycle
(e.g, BN254/Grumpkin). That second curve is folded back in the
first curve. {\bf Decider cost:} \texttt{sonobe} implements
the NOVA + cyclefold. The final decider has to handle
non-native field operations over BN254 (it first uses KZG commitment
to pass values between two curves, but needs to evaluate non-native
field operations of the poly over BN254, and then it eneds to
encode the R1CS group multiplication circuit (on Grumpin) over
BN254 as well. The two operations result in over 11 million
constraints for $0.5 million$ single-step circuit (hence
the overhead of non-native is big, about 10 million constraints).
In Goblin Plonk\footnote{\url{https://hackmd.io/@aztec-network/B19AA8812}}
non-native field operations are encoded as
high-degree (over 30 columns) custom-gates (zk-VM machine),
so the prover's number of constraints is cut down (but 
it does not cut down decider cost). Also it looks like
it is not on one-curve (two proofs? but need to verify).
Reverie \cite{Reverie} further improves CycleFold approach
by mixing it with the ProtoGalaxy folding scheme (multiple instances)
and gives detailed estimate of the cost. Its goal is to reduce
the inter-prover communication, and it is estimated to cut
the cost to around 4000 field elements (around 128kb), which
is used for peer-2-peer network application scenarios (Etherum Consensus).
In Mangrove \cite{Mangrove}, the folding and accumulation scheme
is further extended to tree structure PCD. It divides a plonkish
system into chunks (and using product of grand product to merge
the proof for copy constraints). Thus it achieves very small
memory print for provers. It includes an extension of
Hab22 protocol in its chunking mechanism. However, it requires
the lookup table size not exceeding the entire circuit size so that 
the lookup table (fragments) can fit in the memory of prover.
Note that its Merkle tree structure of all commitments allows to
provide Commit-and-Prove (since original Nova and other schemes
fold witness commitments, hence not easy to take them out).
Recently there are directions in further reducing the
recursive overhead by removing the non-native arithmetics caused
by the homomorphic commitment.
One related needs to mention for Mangrove is Gemini \cite{Gemini},
which also accomplishes space-saving snark via streaming
polynomial commitment and inner product argument prover algorithms
(thus achieving space efficiency). Since Gemini is still
a general purpose working for a large circuit, Mangrove exhibits
better conrete performance via folding compared with Gemini (about
one magnitude).
In SCRIBE \cite{scribe}, streaming is also applied
to sumcheck based prover, where it is shown a magnitude improvement
of speed improvement, and log size RAM access needed (mainly
for a RAM memory for pre-computing evlauation of multi-variate
Lagrange basis polynomials). It looks like
an interesting direction to cut the RAM usage for our decider stage.
In LatticeFold \cite{LatticeFold}, the homomorphic commitment
of witness is replaced by Ajtai's commitment. It 
enjoys the benefits of small field and post-quantum 
security.  The main efforts aims at reducing the increasing norm 
of such commitments, and the depth of folding is limited
to constant round, thus the general PCD tree structure should be used.
NeutronNova \cite{NeutronNova} further applies folding to
the sumcheck protocol via Lagrange interpolation.
In Neo \cite{Neo}, an alternative scheme based on Ajtai's commitment.
It uses sum-check protocol in verifier circuit to prove
that elements are within norm bound and enjoys a claimed
10x speed over HyperNova.

{\bf Comments:} The currently known best folding cost for
Hab22 based lookup argument is ProtoStar and its derivatives
such as ProtoGalaxy and Mangrove. The folding cost seems minimal,
but in the final decider the cost is {\em linear} to the lookup table.
Even if the lookup table size is the same of total queries,
using a SNARK e.g. Plonk will incur more prover cost than
the commit-and-prove approach (e.g., 22MSM for Plonk). It is
possible to mix ProtoStar decider SNARK with the right half of the
CQ protocol, but essentially it still resolves to the 
commit-and-prove approach.

{\bf Comments:} We extended Sonobe \cite{Sonobe}, which instantiates
CycleFold + Nova, exploiting half-pairing Bn254/Grumpkin curves.
Our extension extends it to SuperNova and similarly added
a $24$k constraints for pairing operation at Grumpkin.
MicroNova \cite{MicroNova} similarly uses half-pairing approach
in the context of sum-check based provers (Spartan \cite{Spartan},
which results in slightly higher proof size (using HyperKZG 
to reduce multi-linear polynomial commitments to batched KZG).
MicroNova adopts a variety of techniques to reduce decider size,
e..g, reducing one MLE evaluation over E2 by spliting it into
multiple sub-matrix evaluation. We could adopt the same technique
by spliting the pairing operation into multiple identical cycles,
thus reducing decider size, which remain as a future direction.


\subsection{Lookups}
Many past lookups .... including CQ.


Latest one: {\bf Lasso} \cite{Lasso} which uses sumcheck protocol
and fast multilinear poly commitment to large matrixes with sparse
non-zero entries. It allows to represent (not materialize) a huge
table (e.g., $[0, 2^{128}]$) by projecting a large table to
multiple small tables (decomposable), essentially it's actually
modeling a multi-linear polynomial functions (where smaller functions
are modeled using small look-up tables). This fits very well with
modeling truth table of all CPU instructions, and is applied in
\cite{Jolt}. Note that the prover cost is still dependent
on the container table: for Sona: $O(N)$ field operation snad
$O(\sqrt(N))$ group operations. But verifier has logarithm complexity.
It is shown to be at least 3 times faster than
plookups.  Recently in $\mu$-seek \cite{useek}, the $\cq$ protocol
is modified with multi-linear polynomial commitment to natively
support super-efficient lookups with zk-SNARKs that use 
multilinear commitments; a commit-and-prove construction is
also proved by connection univariate and multi-variate polynomial
commitments. The verifier complexity is logarithmic.


\subsection{Regex Result}
It is well known \cite{AlgBook}[Chapter 5.4] that a regex
can be translated to a graph of $|3m|$ where $m$ is the
regex length. Then to decide if a string of size $n$
belongs to the NFA can be implemented by a BFS search,
thus resulting in $O(mn)$ complexity.

\subsection{Other Regex Efforts}
1. TO READ: ZK-email \url{https://github.com/zkemail/zk-regex}:
According to Reef it converts to DFA and the constraint size
is proportional to DFA size.

2. TO READ: halo2-regex: \url{https://github.com/zkemail/halo2-regex}.
It is using look-ups, but table is inside circuit (cannot exceed
circuit size). Also do not handle large collections of regex.

\subsection{zk-EVM}
Compared with general purpose zero knowledge virtual machines 
such as zk-EVM \cite{polygonevm}, our approach is lighter weight.

In Polygon zk-EVM \cite{polygonevm}, 
a general purpose microprocessor is supported.
Lookups (plookup) is used link state machine traces (lookup up
result produced in another state machine trace).
A RAM machine state is maintained as a Merkle tree.
In some sense, retrieving bag and ternary logic computation
resembles executing micro-instructions on a zk-EVM, however,
as our approach is customized, and it certainly is more
expensive executing it on a general purpose zk-EVM microprocessor.
In Cairo \cite{Cairo}, a general purpose algebraic RISC architecture
is supported. It simulates read-and-write memory using
a memory trace expressed as a read-only dictionary structure
(essentially involving range proofs for proving ascending order
of memory traces and permutation arguments for relating logs
and memory traces). Apparently such a RAM model is more expensive
than our simple mode of passing data between ``big" step-circuits.
Also decoding general machine instructions is more expensive
than a customized model tailored for the specific application domain.
We do want to point out that, using zk-EVM can nicely solve
the problem of decoding an executable file and finding its
entry point, which is approximated by our work. In this case,
how to integrate zk-EVM and customized zk-SNARK circuits would
be an interesting follow-up work.
Recently Nexus \cite{Nexus} further exploits
parallelizable Hypernova \cite{HyperNova} 
and distributed computing to maximize the computing power
of provers. However, customized protocol still shows performance
advantage for specific problems. For instance, its whitepaper
reports that
Nexus VM costs about $30k$ constraints per step, and 
implementing SHA256 usin the general purpose VM is still
about $10^3$ slower than a specialized circuit.

\subsection{RAM}
	We use a global commitment and leverage cyclefold to pass
states between steps in folding schemes. An alternative way
is to use a global RAM as in many variaations of zk-EVM.
General RAM is very expensive, e.g., using
routing network, RSA accumulator, Merkle trees and lookups.
See \cite{AIPRam} for a summary of related work.
Our model is similar to the ``permanent memory" category
in \cite{AIPRam} where there are two commitments to a region of
memory and the prover needs to prove that a program, given the initial
memory setting, leads to the final memory state as committed in the
second commitment. In the latest work, which much improved
the earlier work on memory transcript validation using B\'{e}zout's identity,
ordered permutation leveraging \cite{Hab22}, its complexity 
is $3N + 41A$ (where $N$ is the capacity of RAM and $A$ is the
number of read/write accesses). In our case $A \approx N$,
the general RAM approach is much more expensive, because there
is a large co-efficient over $A$.\footnote{Interestingly, the 
use of B\'{e}zout's identity and the use of CP-Snark in \cite{AIPRam}
is very similar to our earlier work \cite{Zkreg}}.
In StackProofs \cite{StackProofs}, a similar approach to
the Mangrove commit-and-fold is presented, for reasoning about
global memory, which is based on \cite{Hab22}, and it has similar
asymptotic compleixty and we estimate it has similar concrete performance.
In Nebula \cite{Nebula}, commit-and-fold is also used to support
offline RAM verification, improving the state of art of Mangrove by
allowing RAM check in IVC proofs (avoiding generating public verifier
challenge over commitments of all witness at all steps), it also
provides an alternative switchboard approach to SuperNova's approach
in supporting non-uniform circuits, although enjoying the same
asymptotic complexity.

	In HEKATON \cite{HEKATON}, a parallel proof scheme is proposed (note:
not folding). It shares a similar observation with Mangrove that
a long product for finger-printing can be split among multiple circuits.
The key idea is to split a circuit into multiple sub-circuits.
The {\em shared wires} between sub-circuits are modeled as a global
ROM. The memory trace of ROM access is modeled using finger-printing (and
commit and prove snark systems). It improves the concrete
performance of ROM (compared with Mirage+) to $13\times$ the size
of ROM. Notice that however, our approach (in our context) for sharing
incurs only $1\times$ of constraints for I/O memory buffer, because
we do not have to reason about timestamp and trace consistency.
{\em SIDENOTE to remove later: compared with Mangrove the HEKATON
approach pays the extra cost for inter-circuit I/O, if the I/O buffer
is large, it affects concrete performance. }

	In Twist and Shout \cite{TwistShout}, the memory checking for
read/write memory is further improved to only 4$T$ and 6$T$ 
commitments for small and large memory (where $T$ is the trace size),
by leveraging one-hot encoding of addresses and incremental updates
for taking advantage of localicty. The technique does not have a
direct application into our framework given its sumcheck technique
results in log-size proof.

	In the very recent work \cite{RAMLookup} memory checking is
accomplished based on CQ lookup. The idea is that a base table
is prepreprocessed using CQ, and $m$ operations over a deviated
table with hamming distance $\delta$ can be achieved in
complexity $O((m+\delta)\log(m+\delta))$. Then preprocessed CQ keys
are generated periodically, achieving amortized memory access
complexity of $O(m\log(m) + \sqrt{mN})$ per batch. Concretely,
for $2^{20}$ RAM, $10^3$ access costs about $10^2$ seconds (thus
$10$ access per second).  This cannot satisfy our demand of efficiency.

   One recent work zkPi \cite{zkPi} depends on random circuit 
   approach Mirage\cite{Mirage}, which also shows the power of
   the random circuit. (Some of the interesting point, the theoretical
   complexity for non-random circuit to prove permutation is nlog(n),
   but rand circ is O(n), and the analysis of collision for soundness, e.g.,
   $n^3/|F|$).

\subsection{Related in Field}
Compare with zk-MiddleBox, Zombie \cite{Zombie23} and Zk-regex \cite{Luo23}
and Reef \cite{Reef}.

A related line is the TLS Oracle problem where given an TLS session (encrypted), one proves the property of the data (e.g., correct generation of
 authentication key etc.). Traditionally this was done via MPC
 e.g., \cite{Deco,Xie24,Janus} and recently done using zkSNARK in
  Origo \cite{Origo}. Our work can be used to prove TLS data
  with rich set of properties, such as redacted emails, or compliance with
  security policies without revealing encrypted data, efficiently.

\section{Conclusion}
